% Part: lambda-calculus
% Chapter: introduction
% Section: reduction

\documentclass[../../../include/open-logic-section]{subfiles}

\begin{document}

\olfileid{lam}{int}{red}
\olsection{د لامبډا ترمونو راکمول}

د لامبډا ترمونو سره څه کولے شو؟ ساده کول ئې۔ که $M$ او $N$ هر ډول
لامبډا ترمونه او $x$ هر متغير وي، نو له دې وروسته چې د~$M$ هغه تړلي
متغيرونه بيا ونوموو چې تر تعويض وروسته د~$N$ له آزادو متغيرونو سره
ټکرېږي، په~$M$ کښې د~$x$ پر ځاے د~$N$ د کېښودلو پايله په
$\Subst{M}{N}{x}$ ښيو۔ مثلاً،
\[
\Subst{(\lambd[w][xxw])}{yyz}{x} = \lambd[w][(yyz)(yyz)w].
\]

\begin{digress}
د تعويض نورې نښې $[N/x]M$، $[x/N]M$ او $M[x/N]$ هم دي۔ پام
ورته وکړئ!
\end{digress}

په وجداني ډول، $(\lambd[x][M])N$ او $\Subst{M}{N}{x}$ يوه معنا
لري؛ لومړے ترم په دويم بدلول \emph{$\beta$-انقباض} بلل کېږي۔
$(\lambd[x][M])N$ ته \emph{د راکمېدو وړ عبارت} (redex) او
$\Subst{M}{N}{x}$ ته ئې \emph{د انقباض پايله} (contractum) وايي۔
په عمومي ډول، که د~$P$ د کوم فرعي ترم په $\beta$-انقباض سره~$P'$
ته تلل شوني وي، وايو چې $P$ په \emph{يو ګام کښې تر~$P'$ پورې
$\beta$-راکميږي}، او $P \redone P'$ ليکو۔ که له~$P$ څخه د يو-ګامي
راکمول د څو ګامونو په وسيله (ښايي هېڅ ګام نۀ وي)~$P'$ ترلاسه کولے
شو، نو وايو $P$ تر~$P'$ پورې \emph{$\beta$-راکميږي}؛ په نښو کښې
$P \red P'$۔ هغه ترم چې نور $\beta$-راکمول ئې ناشوني وي،
\emph{$\beta$-ناراکمېدونکے} يا \emph{$\beta$-عادي} بلل کېږي۔ چې
سياق روښانه وي، د ``$\beta$-راکميږي'' پر ځاے يوازې ``راکميږي'' او
داسې نور وايو۔

څو بېلګې وګورو۔
\begin{enumerate}
\item لرو:
\begin{align*}
(\lambd[x][xxy]) \lambd[z][z] & \redone (\lambd[z][z])(\lambd[z][z]) y \\
& \redone (\lambd[z][z]) y \\
& \redone y.
\end{align*}
\item د يوه ترم ``ساده کول'' هغه لا پېچلے هم کولے شي:
\begin{align*}
(\lambd[x][xxy])(\lambd[x][xxy]) & \redone (\lambd[x][xxy])(\lambd[x][xxy])y \\
& \redone (\lambd[x][xxy])(\lambd[x][xxy])yy \\
& \redone \dots
\end{align*}
\item ترم بې له بدلون څخه هم پرېښودے شي:
\[
(\lambd[x][xx])(\lambd[x][xx]) \redone (\lambd[x][xx])(\lambd[x][xx]).
\]
\item ځينې ترمونه په څو لارو راکمېدے شي؛ مثلاً،
\[
(\lambd[x][(\lambd[y][yx]) z)] v \redone (\lambd[y][yv]) z
\]
د تر ټولو بهرني تطبيق په انقباض؛ او
\[
(\lambd[x][(\lambd[y][yx]) z)] v \redone (\lambd[x][zx]) v
\]
د تر ټولو دننني تطبيق په انقباض۔ خو پام وکړئ چې په دې حالت کښې
دواړه ترمونه وروسته هماغه يو ترم~$zv$ ته راکميږي۔
\end{enumerate}

په وروستۍ بېلګه کښې دغه ګډه وروستۍ پايله تصادف نۀ دے، بلکې د
لامبډا حساب يو ژور او مهم خاصيت ښيي چې د ``چرچ--روسر خاصيت'' په
نامه يادېږي۔

\end{document}
